\SourceLeaf{450}%
\par\Needspace{9\baselineskip}\begin{center}\textbf{REMARKS.}\end{center}
\noindent\hspace*{1.2em}The communications reproduced on p.\ 397 of this volume, addressed to \textsc{Olbers} and dated 9\ October 1825, show that \textsc{Gauss} began at that time to commit to paper the investigations that form the subject of his \textit{Disquisitiones generales circa superficies curvas} of 1827. The fragment printed above, \textit{Neue allgemeine Untersuchungen über die krummen Flächen}, differs from the \textit{Disquisitiones} in its narrower scope as well as in its method of treatment and the arrangement of the propositions. There \textsc{Gauss} assumes that the rectangular coordinates $x,y,z$ of a point on the surface are represented as functions of any two independent variables $p$ and $q$, whereas here he chooses the geodesic coordinates $s$ and $\theta$ as new variables. Here he first proves that the sum of the three angles of a triangle formed by shortest lines on an arbitrary curved surface differs from $180^{0}$ by the area of the triangle corresponding to it under mapping by parallel normals onto the auxiliary sphere of radius $1$. From this he derives by simple geometric considerations the fundamental theorem that »in the transfer of surfaces by development, the measure of curvature at each point remains unchanged«. In the other work he first shows in \S\ 12 that the measure of curvature can be expressed solely in terms of the three quantities $E,F,G$ and their derivatives with respect to $p$ and $q$, from which the theorem on the preservation of the measure of curvature follows as a corollary. Only much later, in \S\ 20, does he prove the theorem on the angle sum of a geodesic triangle, quite independently of that result.
\begin{flushright}\textsc{Stäckel}.\end{flushright}
